% Lösungen Einheit 1 von 3 – Satz und Hypotenuse (zusammenbau v0.1)
\einheitenkopf{Einheit 1 von 3 · Satz und Hypotenuse}

\erg{12}{a) ja – rechter Winkel bei C}
%% TODO Lösung der Erklärzeile 13g) fehlt in der Bank
\erg{13}{a) Hypotenuse: $41$ cm \quad b) $c^2 = 900 + 1\,600 = 2\,500$, $c = 50$ cm \quad c) $c^2 = 81 + 1\,600 = 1\,681$, $c = 41$ m \quad d) $c^2 = 324 + 576 = 900$, $c = 30$ cm \quad e) $c^2 = 121 + 3\,600 = 3\,721$, $c = 61$ mm \quad f) $c^2 = 1\,600 + 1\,764 = 3\,364$, $c = 58$ m \quad g) \ldots \quad h) $28^2 + 45^2 = 2\,809$, Seite $= 53$ cm \quad i) $c^2 = 97$, $c \approx 9{,}8$ cm \quad j) $c^2 = 42{,}25$, $c = 6{,}5$ m \quad k) $a = 120$ cm; $c^2 = 16\,900$, $c = 130$ cm \quad l) m² + n² = k² \quad m) Kreuz bei der zweiten Aussage: Das Quadrat über der Hypotenuse ist so groß wie die zwei Quadrate über den Katheten zusammen. \quad n) $g = \sqrt{e^2 + f^2}$ (Wurzel aus der Summe der Kathetenquadrate) \quad o) $45^2 + 24^2 = 2\,601$, Diagonale $= 51$ cm \quad p) $60^2 + 11^2 = 3\,721$, Weg $= 61$ m \quad q) schräges Stück $\sqrt{16{,}81} = 4{,}1$ m; Leiter $= 4{,}1 + 1 = 5{,}1$ m \quad r) $AB^2 = 14^2 + 39^2 = 1\,717$, $AB \approx 41{,}4$ m \quad s) $\sqrt{240^2 + 18^2} = \sqrt{57\,924} \approx 240{,}7$ cm}
\erg{14}{a) $4 + 9 = 13$ Kästchen}
\erg{15}{a) $c^2 = 14{,}9$, $c \approx 3{,}9$ m (auf eine Stelle wie die Angaben)}
\erg{16}{a) Wurzel vergessen. Richtig: $c = \sqrt{4\,225} = 65$ cm \quad b) Das Hypotenusenquadrat ist die Summe beider Kathetenquadrate, also größer als jedes einzelne; darum ist auch die Hypotenuse länger als jede Kathete. \quad c) Diagonale $= \sqrt{16\,165} \approx 127{,}1$ cm; $127{,}1 : 2{,}54 \approx 50$ Zoll – ja. \quad d) Wand senkrecht, Boden waagerecht, rechter Winkel zwischen Wand und Boden; die Leiter ist die Hypotenuse H.}

16d)

\dreieckrw{1.88}{5}{$1{,}5$ m}{}{H}

